\documentclass[11pt]{article}
\usepackage[utf8]{inputenc}
\usepackage{amsmath,amssymb}
\usepackage{graphicx}
\usepackage{booktabs}
\usepackage{authblk}
\usepackage[margin=1in]{geometry}
\usepackage[hidelinks]{hyperref}
\usepackage{xurl}
\usepackage{setspace}

\title{Magnetism in RealQM:\\
Currents, the Nuclear Electron, and the Spin Residue}
\author[1]{Claes Johnson}
\affil[1]{KTH Royal Institute of Technology, SE-100 44 Stockholm, Sweden \\
\texttt{claesjohnson@gmail.com}}
\date{\today \\[0.4em]
\small\textit{Text and argument by the author; numerical computations with AI assistance (Claude, Anthropic).}}

\begin{document}
\maketitle

\begin{abstract}
RealQM reformulates quantum mechanics as multiphase continuum mechanics in real three-dimensional space:
electrons are real-valued charge densities on non-overlapping domains with free boundaries, ground states
minimising the Coulomb energy. In that form the theory is deliberately real, and produces no electric
current --- a real wavefunction has $\mathbf J=\frac{q}{m}\,\mathrm{Im}(\psi^*\nabla\psi)\equiv 0$ --- so
magnetism is absent by construction. We develop magnetism as a minimal, current-bearing extension and
follow it, honestly, to where it stops. Four results and one boundary. \emph{(1)}~A complex,
unitary, time-dependent RealQM on the same domains carries a well-defined charge current; the Bernoulli
free boundary automatically conserves charge per domain, so nothing new must be assumed, and the Madelung
hydrodynamic reading becomes a physical flow in real 3D space for any number of electrons. \emph{(2)}~A
circulating charge density carries a magnetic moment $\mu_z=-m\mu_B$, quantised by the winding number
(gyromagnetic $g=1$), and reacts correctly to a field: a numerical solver conserves norm and reproduces
the Zeeman splitting $2\mu_B B$ to four digits. \emph{(3)}~The electron confined in a nucleus is computed
as a \emph{flat, real} density, hence currentless, hence carries \emph{no} moment --- so nuclear moments
come out at nuclear-magneton rather than Bohr-magneton scale, dissolving the classical magnetic-moment
objection to nuclear electrons; the same flatness that makes the electron's mass irrelevant makes its
moment vanish. \emph{(4)}~A single-electron spinor structure delivers $g=2$: from
$(\boldsymbol\sigma\cdot\mathbf D)^2=\mathbf D^2-q\,\boldsymbol\sigma\cdot\mathbf B$ (verified
numerically) the $g=2$ term \emph{emerges}, and an $L=0$ electron splits into two levels with no middle
--- Stern--Gerlach --- non-relativistically. \emph{The boundary:} the many-electron extension fails.
A closed shell comes out paramagnetic, not diamagnetic; forcing the spins antiparallel geometrically
costs energy (a spinor domain wall), so geometry favours parallel. Closed-shell diamagnetism, and
ferromagnetism, require an exchange term that geometric non-overlap does not supply. The precise
statement earned is that RealQM's non-overlap reproduces the \emph{spatial} consequences of antisymmetry
but not its \emph{spin} consequences: single-particle magnetism is delivered; collective magnetism is not.
\end{abstract}

\medskip
\noindent\textbf{Keywords:} RealQM; magnetism; charge current; magnetic moment; nuclear electron; $g=2$;
Stern--Gerlach; exchange; closed shell; spinor.

\doublespacing

\section{Introduction}
\label{sec:intro}

RealQM~\cite{RealQM} represents the $N$ electrons of an atom or molecule by real-valued one-electron
functions $\psi_i:\mathbb R^3\to\mathbb R$ on non-overlapping domains $\Omega_i$ with free (Bernoulli)
boundaries, the ground state minimising the Coulomb energy
\begin{equation}
E=\sum_i\tfrac12\!\int_{\Omega_i}|\nabla\psi_i|^2
+\tfrac12\!\iint\frac{\rho(\mathbf x)\rho(\mathbf y)}{|\mathbf x-\mathbf y|}\,d\mathbf x\,d\mathbf y,
\qquad \rho=\sum_i q_i\psi_i^2,
\label{eq:E}
\end{equation}
found by imaginary-time (parabolic) relaxation. This form is validated across chemistry, but it is
deliberately real, and a real wavefunction carries no current: with $\psi_i$ real,
$\mathrm{Im}(\psi_i^*\nabla\psi_i)=0$. Magnetism --- radiation, moments, susceptibility --- is charge in
motion, and is therefore absent from the base theory by construction.

This paper builds magnetism as a minimal extension and reports, without embellishment, both what it
delivers and where it stops. Sections~\ref{sec:complex}--\ref{sec:spinor} give four results: the
current-bearing formulation, orbital magnetism and its field response, the vanishing moment of the
nuclear electron, and the single-electron $g=2$. Section~\ref{sec:limit} gives the boundary: the
many-electron (closed-shell) extension fails, and we state precisely why. Throughout we set $\hbar=1$ in
formulae; masses $m_i$ and charges $q_i=\pm1$ are kept explicit. Numerical claims are from small,
self-contained solvers documented with the source.

\section{A current-bearing RealQM}
\label{sec:complex}

Let each electron function now be complex and time-dependent, $\psi_i:\Omega_i\times\mathbb R\to\mathbb C$,
with charge density and current
\begin{equation}
\rho_i=q_i|\psi_i|^2,\qquad
\mathbf J_i=\frac{q_i}{m_i}\,\mathrm{Im}(\psi_i^*\nabla\psi_i)=\rho_i\mathbf v_i,\quad
\mathbf v_i=\frac{\nabla S_i}{m_i},
\label{eq:J}
\end{equation}
the last form (writing $\psi_i=R_i e^{iS_i}$) being the Madelung reading: current is charge density times
a velocity set by the phase gradient. The evolution is unitary on each domain,
\begin{equation}
i\,\partial_t\psi_i=\mathcal H_i\psi_i,\qquad
\mathcal H_i=-\frac{1}{2m_i}\nabla^2+V_i,
\label{eq:tdse}
\end{equation}
with $V_i$ the same real (kernel plus inter-domain Hartree) potential as in the ground-state problem, and
the free boundary carrying the \emph{same} Bernoulli condition $\partial_n\psi_i=0$ on $\Gamma_i$.

\paragraph{The free boundary conserves charge.} From \eqref{eq:tdse} with real $V_i$ follows continuity,
$\partial_t\rho_i+\nabla\!\cdot\!\mathbf J_i=0$. The rate of change of charge on $\Omega_i$ is the boundary
flux $-\oint_{\Gamma_i}\mathbf J_i\!\cdot\!\mathbf n$, and the normal current is
$\mathbf J_i\!\cdot\!\mathbf n=(q_i/m_i)\,\mathrm{Im}(\psi_i^*\partial_n\psi_i)$, which the Bernoulli
condition $\partial_n\psi_i=0$ sets to zero. Hence $\dot Q_i=0$: \emph{each electron keeps its unit charge
under the unitary evolution}, because the boundary that defines the RealQM domains is also a no-leak
condition. The static free boundary is exactly what the dynamics needs.

\paragraph{Madelung, physical in 3D.} Equations \eqref{eq:J} are the standard quantum probability current
and the Madelung (1927) hydrodynamic reading. In standard quantum mechanics this reading is physical only
for one particle; for $N$ electrons the fluid lives on $3N$-dimensional configuration space. Because
RealQM keeps every $\psi_i$ in real $\mathbb R^3$, each $(\rho_i,\mathbf v_i,\mathbf J_i)$ is a genuine
three-dimensional density, velocity and current --- one physical charge fluid per electron. RealQM thus
gives the Madelung flow physical existence in real space for any $N$ (with the caveat that this is
RealQM's own fluid, non-overlap standing in for antisymmetry, not a projection of the standard
configuration-space object).

\paragraph{Implementation.} The dynamics is implemented in Cartesian form $\psi_i=u_i+iv_i$, evolved by
the coupled leapfrog $\partial_t u_i=\mathcal H_i v_i$, $\partial_t v_i=-\mathcal H_i u_i$ (Visscher
staggered, norm-preserving to $O(\Delta t^2)$). The polar form $R_i e^{iS_i}$ is avoided: it is singular
at nodes ($S_i$ undefined, $\mathbf v_i$ divergent where $R_i\to0$) --- exactly the axis of a winding
state --- whereas $(u_i,v_i)$ is smooth through nodes and linear.

\section{Orbital magnetism and its response to a field}
\label{sec:orbital}

A stationary state with a spatial phase winding, $\psi=R(\mathbf x)e^{im\phi}$, carries an azimuthal
current and a magnetic moment
\begin{equation}
\boldsymbol\mu=\tfrac12\!\int\mathbf r\times\mathbf J\,d^3x,\qquad \mu_z=-m\,\mu_B,
\label{eq:mu}
\end{equation}
quantised in the winding number and \emph{independent of the density's shape}: the moment-arm ($\propto s$)
and the winding velocity ($\propto 1/s$) cancel, which is the gyromagnetic relation
$\boldsymbol\mu=(q/2m)\mathbf L$ with $L_z=m$ fixed topologically by single-valuedness of $e^{im\phi}$.
This has gyromagnetic factor $g=1$.

\paragraph{Numerical validation.} A Cartesian $(u,v)$ leapfrog for one electron in a 2D harmonic well,
initialised in the $m=1$ winding state $\psi=(x+iy)e^{-r^2/2}$, conserves norm to $1.000000$ and holds
$\mu_z/\mu_B=-1.00$ (grid value $-0.997$) stable over $4000$ steps. The circulating density is genuine.

\paragraph{Reaction to a magnetic field.} Adding the field by minimal coupling
$\mathbf p\to\mathbf p-q\mathbf A$ (uniform $B$ along $z$, $\mathbf A=\tfrac{B}{2}(-y,x)$) and evaluating
the energy of the $m=\pm1$ states gives the Zeeman response (Table~\ref{tab:zeeman}): the circulating
states shift linearly, splitting $E(+1)-E(-1)=2\mu_B B$ to four digits, while the non-circulating $m=0$
state barely moves (weak diamagnetism only). So a RealQM circulating charge density reacts to a field
exactly as a moment $\pm\mu_B$ should --- ordinary orbital magnetism, from a charge density in real space,
no spin invoked.

\begin{table}[h]
\centering\small
\begin{tabular}{cccc}
\toprule
$B$ & $E(m{=}0)$ & $E(+1)-E(-1)$ & $2\mu_B B$ \\
\midrule
$0.05$ & $0.999$ & $+0.0499$ & $0.0500$ \\
$0.10$ & $1.000$ & $+0.0999$ & $0.1000$ \\
$0.20$ & $1.004$ & $+0.1998$ & $0.2000$ \\
\bottomrule
\end{tabular}
\caption{Zeeman response of the circulating density (minimal coupling). The $m=\pm1$ splitting matches
$2\mu_B B$; the $m=0$ (non-circulating) state is inert to linear order.}
\label{tab:zeeman}
\end{table}

\section{The nuclear electron carries no moment}
\label{sec:nuclear}

The classical objection that retired the proton--electron nucleus in 1932 was magnetic: an electron
confined inside a nucleus should carry a moment of order $\mu_B=1836\,\mu_N$, whereas nuclear moments (the
neutron's $-1.913\,\mu_N$ included) are of order $\mu_N$ --- a thousandfold too small. In a charge-density
theory the moment is the current's \eqref{eq:mu}, and a real density has $\mathbf J=0$. RealNucleus
computes the deuteron's confined electron as a \emph{flat, real, node-less} density~\cite{RealNucleus};
it is the $m=0$, currentless case, and carries \emph{no} magnetic moment. Nuclear moments are then set by
the protons and come out at $\mu_N$ scale. The thousandfold overshoot does not occur, because in a
charge-density theory there is no intrinsic moment to carry --- only the current, and the flat electron's
current is zero.

\paragraph{Energetics, and one flatness for two objections.} A winding at confinement scale $\sigma$ costs
$\Delta T=\tfrac{\hbar^2}{2m_e}m^2\langle1/s^2\rangle\approx\mathrm{Ry}\,(a_0/\sigma)^2$: about one Rydberg
($27$~eV) for a free electron ($\sigma\sim a_0$) but $\sim$GeV for a femtometre-caged one --- thousands of
times the deuteron binding, so the caged electron cannot afford a winding and settles at $m=0$. A radial
relaxation confirms this self-consistently. And the point is the \emph{same} flatness that was introduced
to make the electron's \emph{mass} irrelevant to nuclear binding (a flat density has vanishing kinetic
gradient): that flatness, being real and currentless, makes the \emph{magnetic moment} vanish too. Two
of the classical objections to nuclear electrons --- confinement energy and magnetic moment --- are
answered by one property of the confined-electron solution.

\section{Single-electron spin: $g=2$ from a spinor structure}
\label{sec:spinor}

Orbital circulation gives $g=1$; the electron's spin moment has $g=2$ (a spin of $\hbar/2$ yields a full
$\mu_B$), and the two-valued Stern--Gerlach response is not orbital at all. This is not a demand for
relativity. By Lévy-Leblond~\cite{LevyLeblond}, $g=2$ follows from a first-order, \emph{non-relativistic}
spinor equation. Equipping the domains with a two-component spinor and the linearised structure
$(\boldsymbol\sigma\cdot\mathbf D)$, $\mathbf D=\mathbf p-q\mathbf A$, the operator identity
\begin{equation}
(\boldsymbol\sigma\cdot\mathbf D)^2=\mathbf D^2\,\mathbb 1-q\,\boldsymbol\sigma\cdot\mathbf B
\label{eq:pauli}
\end{equation}
holds (verified numerically to grid accuracy, residual $\sim10^{-3}$), so the $g=2$ Zeeman term
$-q\,\boldsymbol\sigma\cdot\mathbf B$ \emph{emerges} from squaring $\boldsymbol\sigma\cdot\mathbf p$ --- it
is not inserted. Consequently an $L=0$ (orbitally inert) electron --- the state a scalar density leaves
undeflected --- splits under the spinor into \emph{two} levels $\pm\mu_B B$, splitting $2\mu_B B$,
$g_s=2$, with \emph{no middle}, against the three-level ($m=-1,0,+1$) orbital pattern. This is
Stern--Gerlach, delivered non-relativistically: spin reaches $\pm\mu_B$ at $\pm\tfrac12$ unit of angular
momentum ($g=2$, no middle) where orbital needs $\pm1$ unit ($g=1$, with middle). The ``2'' is the SU(2)
double cover of the rotation group --- geometric, not relativistic.

\section{The boundary: collective magnetism is not supplied}
\label{sec:limit}

Here the extension stops, and it is important to say so plainly. ``Spin'' bundles two logically distinct
things: \emph{statistics-spin} (the fermionic exclusion --- two electrons per orbital, the periodic
table), which RealQM realises geometrically by non-overlap with no spin variable; and \emph{magnetic-spin}
(the two-valued moment of \S\ref{sec:spinor}). The single-electron magnetic-spin is delivered. The
many-electron case is not.

\paragraph{The closed-shell test fails.} A filled shell must be diamagnetic --- two electrons pairing to
zero net moment. We tested the spin sector of the naive extension: two electrons, each with an independent
spinor, in a field. The ground state is the aligned pair with net moment $2\mu_B$ --- \emph{paramagnetic},
not diamagnetic. Nothing in geometric non-overlap forces the two spins antiparallel.

\paragraph{Geometry does not rescue it.} One might hope continuity of the spinor field across the shared
domain boundary would force the spins antiparallel. It does the opposite. Forcing the joined field
antiparallel means a spinor \emph{domain wall}, whose kinetic cost is
$T_{\rm spin}=\tfrac18\!\int(\theta')^2\approx\pi^2/24w>0$, while the parallel configuration costs nothing:
continuity \emph{favours parallel} (paramagnetic). And in the actual RealQM structure --- separate
functions on non-overlapping domains --- the spinors are simply uncoupled and degenerate, aligning with
any field. Either way geometry supplies no antiparallel pairing. Obtaining diamagnetism requires an
antiparallel-pairing (exchange) term $J\,\mathbf S_1\!\cdot\!\mathbf S_2$, which is a spin-statistics
ingredient non-overlap does not carry. The same missing exchange is what drives \emph{collective}
magnetism generally: ferromagnetism is the spontaneous alignment of atomic moments by exchange, and it too
is absent.

\paragraph{The precise statement.} Geometric non-overlap reproduces the \emph{spatial} consequences of
antisymmetry (exclusion, shells, the whole of chemistry) but not its \emph{spin} consequences (singlet
pairing, exchange splitting, closed-shell diamagnetism, ferromagnetism). Single-particle magnetism is
delivered; collective, exchange-driven magnetism is a genuine limitation, not a gap awaiting a boundary
condition.

\section{Interpretive frame: magnetism without relativity}
\label{sec:frame}

The results sit naturally in a reading of RealQM in which an electron is not a standalone object but a
\emph{domain of a configuration}, its size fixed by Coulomb balance within the configuration rather than
by any intrinsic scale~\cite{Charges}. On that reading the single-atom magnetic moment is a configuration
property --- $\mu_B$ when the electron is loosely confined (a winding it can afford), none when tightly
caged --- and no rest-energy or Compton scale is invoked anywhere; even $g=2$ is the geometric SU(2)
factor, not a relativistic one. This frame is not needed for the computational results of
\S\S\ref{sec:complex}--\ref{sec:limit}, which stand on their own; it is offered as the picture that makes
them coherent, and as the reason the one true residue --- collective magnetism --- is a matter of missing
\emph{exchange}, not of missing relativity.

\section{Conclusion}
\label{sec:concl}

Magnetism enters RealQM as a minimal current-bearing extension: complex, unitary dynamics on the same
domains, with the free boundary conserving charge and the Madelung flow made physical in real 3D. From it,
orbital magnetism is delivered and reacts correctly to a field; the flat confined electron carries no
moment, dissolving the magnetic-moment objection to nuclear electrons by the same flatness that removes
its mass; and a single-electron spinor structure yields $g=2$ and Stern--Gerlach non-relativistically. The
extension then reaches a firm boundary: closed-shell diamagnetism and ferromagnetism require an exchange
term that geometric non-overlap does not supply, because non-overlap carries the spatial but not the spin
consequences of antisymmetry. What RealQM delivers is single-particle magnetism, honestly and without
relativity; what it does not deliver is the collective, exchange-driven magnetism that binds many moments
together --- and locating that boundary precisely is itself a result.

\begin{thebibliography}{9}
\bibitem{RealQM} C.~Johnson, \emph{Real Quantum Mechanics as Multiphase 3D Continuum Mechanics},
manuscript, 2026; \url{https://claes542.github.io/RealMolecule/gallery.html}.
\bibitem{RealNucleus} C.~Johnson, \emph{RealNucleus: Nuclear Binding as Dual Coulomb Confinement, without
a Strong or Weak Force}, 2026.
\bibitem{Charges} C.~Johnson, \emph{Elementary Charges as Configurations, not Objects: Why RealQM Needs No
Relativity}, 2026.
\bibitem{LevyLeblond} J.-M.~Lévy-Leblond, \emph{Nonrelativistic particles and wave equations},
Commun. Math. Phys. \textbf{6}, 286 (1967).
\end{thebibliography}

\end{document}
